% Option-A replacement for the Introduction section.
% This fragment is intentionally conservative until the official A1/A2/A3 campaign is complete.

\section{Introduction}\label{sec:introduction}

The expansion of blockchain systems beyond single-chain architectures has created a multi-chain ecosystem in which decentralized applications transfer assets and messages across otherwise independent ledgers~\citep{belchior2021survey,augusto2024sok}. Cross-chain bridges are a critical part of this infrastructure, but they also enlarge the security boundary: a transfer may depend on source-chain contracts, off-chain or validator-mediated message handling, and destination-chain contracts whose states evolve asynchronously.

Historical bridge incidents illustrate the practical impact of failures across these domains. Publicly documented attacks include failures in validator or authorization logic, incorrect message verification, initialization errors, forged or inconsistent deposit evidence, and state-transition bugs that become exploitable only when source, relay, and destination behavior are considered together~\citep{lee2023sok,duan2023attacks,liao2024smartaxe}. Reported loss figures vary with valuation timestamp, recovery status, and whether a source reports gross outflow or net unrecovered loss; we therefore use them only as contextual indicators rather than as experimental measurements.

Existing defenses address complementary parts of the problem. Cross-chain static analyzers reason about source-level data and control flow~\citep{liao2024smartaxe,zhou2025bridgeguard}; monitoring systems check observed bridge executions or business-logic conformance~\citep{zhang2022xscope,eshghie2024highguard,augusto2025xchainwatcher}; and graph-based approaches detect suspicious cross-chain transaction behavior~\citep{wu2025safeguarding,lin2025bridgeshield}. Dynamic testing has also moved directly into the cross-chain setting. BridgeFuzz exercises bridge deployments with on-chain and off-chain relayer scope~\citep{winkler2026bridgefuzz}, while IntentFuzz targets intent-based bridges using automatically recovered protocol roles, invariant-oriented execution, and LLM-assisted input construction when deterministic generation is insufficient~\citep{augusto2026intentfuzz}. Consequently, BridgeSentry does not claim novelty for cross-chain fuzzing, relayer-aware execution, or LLM-assisted fuzzing in isolation.

The narrower question studied in this work is how heterogeneous bridge artifacts can be translated into explicit cross-domain security properties and execution guidance while keeping historical knowledge separate from the evidence used to validate an exploit trace. This separation matters because target-aware reconstruction and unseen-vulnerability discovery are different empirical tasks. A system that retrieves the description of a known incident and reproduces its target predicate may be useful for testing a reconstruction harness, but that success does not establish generalization to an incident, bridge family, or vulnerability class that was excluded from the guidance corpus.

We present \textbf{BridgeSentry}, a semantic-guided framework for cross-chain exploit reconstruction and validation. BridgeSentry organizes analysis around an explicit source--relay--destination representation and uses three coordinated stages:

\begin{enumerate}
    \item \textbf{Hybrid semantic extraction and candidate invariant synthesis.} Bridge artifacts are processed using structured program analysis when available, a conservative parser fallback, and optional LLM enrichment. The resulting entities, guards, asset/message flows, and dependencies are compiled into an Atomic Transfer Graph (ATG) together with candidate protocol properties. Generated candidates are not treated as validated specifications by default.

    \item \textbf{Provenance-aware historical guidance.} A structured corpus of documented bridge incidents is retrieved to construct threat-model-constrained adversarial action sequences and semantic waypoints. Every generated scenario retains retrieval provenance so that full-knowledge reconstruction can be separated from incident-, bridge-family-, vulnerability-class-, and temporal-holdout regimes.

    \item \textbf{Dual-EVM execution with fail-closed evidence validation.} A source EVM and destination EVM interact through an explicit relay model. Scenarios are used as seeds and guidance, whereas exploit evidence is required to come from observed execution state, transaction outcomes, logs, storage changes, relay state, and paired patched controls. Missing evidence is treated as unevaluable rather than implicitly successful.
\end{enumerate}

This framing deliberately separates three questions that are often conflated: whether a system can recover bridge semantics, whether historical knowledge can guide reconstruction, and whether that guidance generalizes when target-specific knowledge is excluded. The retained legacy experiments in the current artifact address only the second question under a target-aware setting. They contain 12 historical reconstruction fixtures---10 EVM reconstructions, one EVM surrogate of a Solana vulnerability, and one scenario-layer-only case---and report 12/12 target-predicate triggers, 11 cases with target-bytecode execution, and a mean fuzzing-only trigger time of 1.66 seconds on successful legacy runs. We do not interpret these quantities as Valid Exploit Rate (VER), precision, or zero-day discovery performance because the legacy corpus is target-leaked, positive-only, and lacks a completed paired-patch and benign-control campaign.

The contribution of this paper is therefore structured around evidence that can be evaluated independently:

\begin{itemize}
    \item \textbf{A cross-domain semantic intermediate representation.} We adapt the Atomic Transfer Graph abstraction~\citep{dubler2025atg} to represent source, relay, and destination dependencies together with canonical per-transfer/per-message security properties for dynamic bridge testing. The ATG is used as an intermediate representation between artifact analysis and execution, not as a secure-by-construction specification.

    \item \textbf{Leakage-aware historical guidance.} We treat the historical incident corpus as both a source of useful attack knowledge and a potential source of target leakage. Retrieval identifiers and corpus versions are retained per run, and effectiveness claims are conditioned on explicit knowledge-exclusion regimes rather than inferred from full-knowledge reconstruction.

    \item \textbf{A fail-closed definition of exploit evidence.} We distinguish an intermediate invariant trigger from a bytecode-backed, causally valid, capability-compliant, reproducible trace that causes material impact and is rejected by a paired patched revision. Unknown or unobserved required evidence makes a trace ineligible for VER.

    \item \textbf{An empirical evaluation protocol that separates semantic fidelity, guidance generalization, exploit validity, and component attribution.} The final evaluation uses labelled semantic data, leakage-controlled retrieval splits, vulnerable/patched pairs, benign and hard-negative controls, and right-censored analysis for stochastic time-to-valid-exploit outcomes. The legacy 12-incident reconstruction study is retained only as historical diagnostic evidence.
\end{itemize}

The rest of the paper is organized as follows. Section~\ref{sec:related} reviews cross-chain security analysis, dynamic testing, and LLM/retrieval-assisted security techniques. Section~\ref{sec:problem} defines the execution model, ATG representation, protocol properties, adversary capabilities, and the distinction between reconstruction and restricted-knowledge generalization. Section~\ref{sec:methodology} presents the BridgeSentry pipeline. Section~\ref{sec:experiments} defines and reports the evaluation. The remaining sections discuss limitations and threats to validity and conclude the paper.
